\chapter{Complex Amplitudes, and an Obstruction}
\label{sec:quantum}
% ===================================================================

\begin{remark}[Status of this chapter]
\label{rmk:ch09status}
This chapter is the least finished thing in this part of the book, and the
reader should treat it accordingly.  An earlier version presented the
complex instance of the decoration as delivering quantum dynamics
outright.  It does not, and the reason it does not is structural rather
than technical --- it will not be repaired by working harder at the same
construction.  What follows states the construction, states the
obstruction, and says what would have to change.  Everything after
Section~\ref{sec:obstruction} is a research programme, not a result.
\end{remark}

\section{The complex instance}

Take $\Semi = \CC$ in Definition~\ref{def:decoration}, with the reversal
condition $\dec^{\dagger}_r(\phi) = \overline{\dec_r(\phi)}$ of
Chapter~\ref{sec:extended-modal}.  Definition~\ref{def:pathweight} then
reads: the amplitude of a path is the product of its step amplitudes, and
the transition amplitude from $P$ to $Q$ is
\[
  \langle Q \mid P \rangle \;=\; \sum_{\gamma : P \to Q} W(\gamma),
\]
with $|\langle Q \mid P\rangle|^2$ the candidate transition probability.
This much is a construction and is not in doubt.  It is also, so far,
only bookkeeping: complex numbers have been attached to steps and
multiplied along paths.

What would make it quantum is \emph{cancellation}.  Two paths $\gamma_1,
\gamma_2$ from $P$ to $Q$ with $W(\gamma_1) = -W(\gamma_2)$ should
annihilate, so that the transition has zero probability despite having two
distinct causal histories.  That is the phenomenon; everything else is
notation.

\section{The obstruction}
\label{sec:obstruction}

Cancellation is exactly what this framework, as built, forbids.

\begin{remark}[Interference versus branching time]
\label{rmk:obstruction}
The equivalence on which this entire development rests is bisimulation,
and bisimulation is a \emph{branching-time} notion: it distinguishes
systems by the shape of their tree of possibilities, not merely by the set
of runs they admit.  Two distinct paths to the same term are, for
bisimulation, two distinct pieces of structure, and no assignment of
weights to steps can make a branching-time equivalence identify them ---
still less make them annul one another.  Interference adds and
\emph{cancels} contributions to a shared outcome, which is a statement
about runs; bisimulation refuses to pool runs in the first
place~\cite[\S13]{CE}.

The two are therefore not merely in tension.  The genuinely quantum
reading pulls toward a \emph{linear-time} semantics, and that is a
reformulation of the framework rather than a relabelling of it.
\end{remark}

The obstruction is worth stating sharply because it is easy to disguise.
One can write down $\sum_\gamma W(\gamma)$ and observe that terms cancel;
this is arithmetic and always available.  What one cannot do, without
changing the semantics, is claim that the cancellation is \emph{observed}
by the theory --- that the resulting system is behaviourally
indistinguishable from one in which the transition never occurs.  Under
bisimulation it is not.  The two paths remain visible.

\begin{remark}[What was previously claimed about density matrices]
An earlier version of this chapter proposed letting the vectorial account
$\vect{A}$ become a positive semidefinite matrix $\rho$, with
$\mathrm{tr}(\rho)=1$ read as the quantum analogue of resource
conservation.  This should not be retained.  A resource account is an
element of an ordered monoid, tracking what may be spent; a density matrix
is a state, tracking what may be observed.  They are not the same kind of
thing, and identifying them conflates the conservation law of
Proposition~\ref{prop:firstlaw} --- which is about value under conversion
--- with the normalisation of a probability distribution.  Whatever the
right quantum statement is, it is not that one.
\end{remark}

\section{What would have to change}

Three routes are open, and we have not chosen among them.

\begin{enumerate}[itemsep=5pt]
  \item \textbf{Move to linear time.}  Replace bisimulation by a
        trace-based equivalence, accept the loss of the branching-time
        distinctions the rest of this book depends on, and recover
        interference honestly.  The cost is high: the ontology of
        Part~\ref{part:sciencebot}, in which bisimulation classes are the
        real things, does not survive the move intact.

  \item \textbf{Quotient late.}  Keep bisimulation as the semantics and
        introduce cancellation as a separate extraction laid on top ---
        the Born rule as an additional map out of the model rather than a
        property of it.  This is the categorical-quantum-mechanics
        pattern, where the interaction structure yields a semiring of
        scalars and the probability rule is applied
        afterwards~\cite{abramsky}.  It is the most conservative route and
        the least explanatory.

  \item \textbf{Find the right equivalence.}  Ask whether there is a
        coarsening of bisimulation that pools exactly the runs quantum
        mechanics pools and no others --- interference as a \emph{quotient
        of branching structure} rather than as an alternative to it.  We
        regard this as the interesting question and have nothing to report
        on it.
\end{enumerate}

\section{Simulation, in the meantime}

The stochastic instance, by contrast, is unobstructed, and it is what the
implementation actually runs.  The classical Gillespie
algorithm~\cite{gillespie1977exact} simulates a continuous-time Markov
chain by computing exit rates from the current state, sampling a waiting
time, and selecting the next rule proportionally to rate.  With
$\Semi = \Real_{\geq 0}$ and a dynamic decoration
(Definition~\ref{def:dyndec}) supplying the rates, this is precisely a
decorated GSLT stepped forward: the base rewrite rules of a language
definition carry the decoration, and the simulator evaluates a Gillespie
step at each state, advancing the term, updating the table, debiting the
account and recording the trace~\cite{StochSim}.

Quantum analogues of the Gillespie algorithm exist for open systems
described by Lindblad master equations, with the exit rates becoming
complex decay parameters~\cite{cao2021quantum}.  Nothing prevents running
that loop over a complex decoration, and the numbers that come out will
look right.  Remark~\ref{rmk:obstruction} is the reason we decline to call
the result a quantum dynamics: the simulation would be sampling a linear-time
object from a branching-time model, and until the mismatch is resolved the
output is a computation in search of a semantics.

% ===================================================================
